\documentclass[12pt,norsk]{beamer}

\usepackage[norsk]{babel}
\usepackage[utf8]{inputenc}
\usepackage{graphicx}

\usetheme{Luebeck}
\usecolortheme{dolphin}

\graphicspath{ {images/} }




%----------------------------------------------------------------------------------------
%	TITLE PAGE
%----------------------------------------------------------------------------------------

\title[Avoiding The Top 10 Software Security Design Flaws]{Avoiding The Top 10 Software Security Design Flaws}
% \subtitle{Chapter 6.1 - 6.2}

\author{Lasse K. Brun \& Haavard R. Olsen \\* (Group 5)} % Your name
\institute[UiB]
{
\medskip
\textit{lkbrun@gmail.com \& haavard.olsen@live.com} \\* % Your email address
\textit{INF226 - Software Security, Fall 2014 }
}
\date{13. November 2014} % Date, can be changed to a custom date



\begin{document}

 

\begin{frame}
\titlepage % Print the title page as the first slide
\end{frame}

%----------------------------------------------------------------------------------------
%	PRESENTATION SLIDES
%----------------------------------------------------------------------------------------

\begin{frame}
	
	\frametitle{Introduction}

	\begin{itemize}
		\item IEEE Computer Society Center for Secure Design.
		\item Shifting focus from finding implementation bugs to identifying common design flaws.
		\item Brought together people from different organizations at a workshop in 2014.
		\item Discussed types of flaws
		\begin{itemize}
			\item they identified in their own internal design reviews.
			\item that were available from external data.
		\end{itemize}
		\item Resulted in the list: Avoiding the top 10 software security design flaws.\cite{t10ssdf}
	\end{itemize}		

\end{frame}


\begin{frame}
	
	\frametitle{Contributors}
	
	\begin{figure}[h]
    \centering
    \includegraphics[scale=0.50]{images/contributors.png}
    
\end{figure}
	
\end{frame}

\begin{frame}
	
	\frametitle{Mission Statement}

	\begin{itemize}
		\item Recognizing software system designs that are likely to compromise.
		\item Designing and building software systems with strong, identifiable security properties.
	\end{itemize}		

\end{frame}

\begin{frame}
	
	\frametitle{Shifting focus}

	\begin{itemize}
		\item Today's focus lies generally on finding implementation bugs, rather than design flaws.
		\item IEEE wants to shift the focus towards design flaws.
		\item A bug is an implementation-level software problem.
		\item A design flaw is a problem at a deeper level, and is the result of a mistake or oversight at the design phase.
		
	\end{itemize}		

\end{frame}

\begin{frame}

	\frametitle{1. Earn or Give, but never assume, Trust}
	
	\begin{itemize}
		\item Think about where to implement security functionality. 
		\item Consider all types of clients.
		\item Consider client input compromised until proven otherwise.
		\item Safer in the long run to assume a system not trustworthy.
		\item Design client side to cope with potential compromise.
		\item Validate user input.
		\item Consider where code will be executed, where data will go and where data entering your system comes from. 
	\end{itemize}

\end{frame}

\begin{frame}

	\frametitle{2. Use an authentication mechanism that cannot be bypassed or tampered with}
	
	\begin{itemize}
		\item One goal of secure design: prevent unauthenticated access.
		\item Prevent changing of identity without re-authenticate.
		\item Authenticating requires one or more factors.
	
		\begin{itemize}
			\item Something you know (Password)
			\item Something you are (Biometric signature)
			\item Something you have (Smartphone)
		\end{itemize}
		
		\item Authenticate machines as well as humans.	
	
	\end{itemize}
	
	
\end{frame}

\begin{frame}

	\frametitle{2. Use an authentication mechanism that cannot be bypassed or tampered with}
	
	\begin{itemize}
		\item Don't use forgeable session tokens.
		\item Use time-tested mechanisms such as Kerberos\cite{kerberos}. 
		\item Specify time limit for the session if user is inactive.
		\item Handle passwords properly!
		\item It's preferable to use on component responsible for authentication.
	\end{itemize}
	
\end{frame}

\begin{frame}

	\frametitle{3. Authorize after you authenticate}

	\begin{itemize}
		\item Determine what actions users can execute on the system.
		\item Authentication may not be sufficient to deciding whether to disallow or allow the user to perform certain actions. 
		\item Authorization depend upon privileges associated with user and context of the request.
		\item Authorization may need to invoke authentication.
		\item A common infrastructure should be responsible for authenticating users.
		
	\end{itemize}		

\end{frame}

\begin{frame}

	\frametitle{4. Strictly separate data and control instructions, and never process control instructions received from untrused sources}
	
	\begin{itemize}
	
		\item Co-mingling data and control instructions in a single entity can lead to injection vulnerabilities.
		\item General problem in all layers of computer systems.
		\item At Lower levels:
	
		\begin{itemize}
			\item Memory-corruption vulnerabilities
			\item May cause attacker controlled modification of control flow
			\item May cause direct execution of data as machine or bytecode instructions
		
		\end{itemize}
	
	\end{itemize}

\end{frame}
\begin{frame}
	\frametitle{4. Strictly separate data and control instructions, and never process control instructions received from untrused sources}
	
	\begin{itemize}
	
		\item At Higher levels:
		\begin{itemize}
			\item Problems occur in context of runtime interpretation of domain-specific or general-purpose programming languages.
			\item If sw assembles string in a parseable language by combining untrusted data with control instructions.
		\end{itemize}
	
		\item You should consider control-flow integrity and separation of control and data as important design goals!
		\item What to consider while designing/using APIs! - Class discussion!
		 
	\end{itemize}

\end{frame}

\begin{frame}
	\frametitle{4. Strictly separate data and control instructions, and never process control instructions received from untrused sources}
	
	\begin{itemize}
		\item Designing systems which relies on transforming data into code:
		\begin{itemize}
			\item Take extra care to validate the data as strictly as possible.
			\item Areas this is concerning is the \textit{\textbf{eval}} function, reflection or query languages.
		
		\end{itemize}
	
	\end{itemize}

\end{frame}

\begin{frame}

	\frametitle{5. Define an approach that ensures all data are explicitly validated}
	
	\begin{itemize}
		\item Ensure that comprehensive data validation actually takes place.
		\item All assumptions about data have been validated when they are being used.
		\item Feasible for security reviewer to efficiently reason about and verify the correctness of data validation.
		\item This verifiability should take into account that code change.
	
	\end{itemize}
	

\end{frame}

\begin{frame}

	\frametitle{5. Define an approach that ensures all data are explicitly validated}
	
	\begin{itemize}
		\item Design and use centralized validation to ensure that all data entering the system is properly validated. 
		\begin{itemize}
			\item Intercept all incoming requests.
			\item Validate all fields before any processing is done.
			\item Consuming complex data formats, must such parsers be designed to cope with malicious input.

		\end{itemize}

		\item Transform data to canonical form.
		\item Use common libraries of validation of primitives. Should be white listing. 			
	
	\end{itemize}
	

\end{frame}

\begin{frame}

	\frametitle{6. Use cryptography correctly}

	\begin{itemize}
		
		\item One of the most important tools for writing secure software.
		\item What cryptography can do:
	
		\begin{itemize}
	
			\item Ensure confidentiality	
			\item Protect from unauthorized modification
			\item Authenticate the source of data
		\end{itemize}	 
	
		\item Extremely hard to get right!
		\item Long list of pitfalls...	
	
	\end{itemize}

\end{frame}

\begin{frame}

	\frametitle{6. Use cryptography correctly}
	
	\begin{itemize}
		
		\item Using your own cryptographic algorithms and implementation.
		\item Misuse of libraries and algorithms.
		\item Poor key management.
		\item Bad randomisation algorithms.
		\item Failure to centralize cryptography.
	\end{itemize}
\end{frame}

\begin{frame}

	\frametitle{7. Identify sensitive data and how they should be handled}
	\begin{itemize}
		\item Identify sensitive data and how to protect it.
		\item Many systems have failed to protect data over the years.
		\item Data sensitivity is context-sensitive:
		\begin{itemize}
			\item Regulations
			\item Company policies
			\item Contractual obligations
			\item User expectations
		\end{itemize}
		\item Crucial to factor all relevant considerations into the design.
		
	\end{itemize}
	

\end{frame}

\begin{frame}

	\frametitle{7. Identify sensitive data and how they should be handled}
	\begin{itemize}

		\item Not all data protection requirements are the same.
		\begin{itemize}
			\item Confidentiality: Financial records and corporate intellectual property.
			\item Availability: Medical records.
			\item Integrity: Spoofing or substituting data to cause a system to misbehave intentionally.
		\end{itemize}
		
		
	\end{itemize}
	

\end{frame}

\begin{frame}

	\frametitle{7. Identify sensitive data and how they should be handled}
	\begin{itemize}

		\item Protect data in transit between system components and between organizations.
		\item Data crossing trust boundaries.
		\item Essential to identify trust boundaries and rectifying them with protection policies.
		
		
		
	\end{itemize}
	

\end{frame}


\begin{frame}

	\frametitle{7. Identify sensitive data and how they should be handled}
	\begin{itemize}

		\item Policy requirements and data sensitivity can change over time.
		\item Regulatory revisiting and revising data protection policies.		
		
	\end{itemize}
	

\end{frame}

\begin{frame}

	\frametitle{8. Always consider the users}
	
	\begin{itemize}
	
		\item Security directly linked to the users.
		\item Balance between good user experience and sufficient security.
		\item Consider what system you are creating as well as who will use it.
		\item Recommendation: "Secure by default".
	\end{itemize}

\end{frame}

\begin{frame}
 
	\frametitle{8. Always consider the users}
	
	\begin{itemize}
	
	\item Alot of work to be done in user-focused security!
	\item Example problem: "Secure" password requirements.
	\item More complex: Simple security for normal users, while giving admins the control they require.
	
	\end{itemize}

\end{frame}

\begin{frame}

	\frametitle{9. Understand how integrating external components changes your attack surface}
	
	\begin{itemize}
		\item Unlikely to develop a system without external components.
		\item When adding functionality, developers often use existing components or external components.
		\item External components are ''not written here".
		\begin{itemize}
			\item Off the shelf components, platforms, applications.
			\item Third party open source.
			\item Widgets or gadgets added to a web project.
			\item Web service.
		\end{itemize}
		\item Can be included as binaries, libraries, source code or simple API's.
	\end{itemize}
\end{frame}

\begin{frame}

	\frametitle{9. Understand how integrating external components changes your attack surface}
	
	\begin{itemize}
		\item Functionality should be achieved by reuse, rather than creating something from scratch.
		\item These have been tested and validated to meet a certain security standard.
		\item By using rather than building means that the software inherits the security weaknesses.
		\item This must be solved, mitigated or accounted for.
	\end{itemize}
	

\end{frame}


\begin{frame}

	\frametitle{9. Understand how integrating external components changes your attack surface}
	
	\begin{itemize}
		\item Allocate time in software development methodology to consider security impact.
		\begin{itemize}
			\item Change threat model? 
			\item Add any attack surface?
			\item Unused features?
			\item Included components include other components?
			\item Source known and trusted? 
			\item Documentation?
		\end{itemize}		 
		\item External components should not be trusted until appropriate security controls have been applied.
	\end{itemize}
	

\end{frame}


\begin{frame}

	\frametitle{9. Understand how integrating external components changes your attack surface}
	
	\begin{itemize}
		\item At the minimum consider the following:
		\begin{itemize}
			\item Only enable functionality you intend to use.
			\item If you can't configure security properties, use another library. 
			\item Follow sources that track or publish security-related information.
			\item Maintain an up to date list of all external components.
			\item Consider all data from external sources tainted  until proven valid.
			\item Design for flexibility.
		\end{itemize}
		
	\end{itemize}
	

\end{frame}

\begin{frame}

	\frametitle{10. Be flexible when considering future changes to objects and actors}
	\begin{itemize}
		\item Security designed for change!
		\item Example: Change in threats or environment.
		\item How to keep your system flexible?
	\end{itemize}
	
	

\end{frame}

\begin{frame}

	\frametitle{10. Be flexible when considering future changes to objects and actors}
	\begin{itemize}
		\item Design for secure updates.
		\item Design for security properties changing over time.
		\item Design with the ability to toggle functionality.
		\item Design for changes to secret objects.
		\item Design for changes in third-party applications.
		\item Design for changes to entitlements.
	\end{itemize}
	
	

\end{frame}

\begin{frame}

	\frametitle{Our opinion on T10SSDF}

	\begin{itemize}

		\item Pros:
		
		\begin{itemize}
			
			\item A lot of valid and good arguments.
			\item A lot of interesting and useful information.		
		\end{itemize}				
				
		\item Cons:
		
		\begin{itemize}
			
			\item Hard to read.
			\item Complicates simple problems and points.
			\item Compresses pages of text into a few sentences.
			
		\end{itemize}
	\end{itemize}

\end{frame}

\begin{frame}
\frametitle{Bad example}

\textit{ `` The challenge to designers and architects lies in creating designs that facilitate secure configuration and use by those interested in doing so, designs that motivate and incentivize secure use among those not particularly interested in software security, and designs that prevent or mitigate abuse from those who intend to weaken or compromise the system. "}

\end{frame}

\begin{frame}

	\frametitle{Questions?}
		

	\begin{figure}[h]
    	\centering
    	\includegraphics[scale=0.45]{images/questionmark.png}
    
	\end{figure}

\end{frame}


\begin{frame}

	\frametitle{References}
		
	\bibliographystyle{plain}
	\bibliography{biblio}	

\end{frame}




%----------------------------------------------------------------------------------------

\end{document} 